\documentclass[a4paper,12pt]{article}
\usepackage[italian]{babel}
\usepackage[utf8]{inputenc}
\usepackage[T1]{fontenc}
\usepackage{geometry}
\usepackage{setspace}
\usepackage{amsmath}
\geometry{margin=2.5cm}
\onehalfspacing

\title{Soluzioni degli Esercizi sulle Subnet}
\author{Prof. Fedeli Massimo}
\date{}

\begin{document}

\maketitle

\section*{Esercizio 1}
Dato \textbf{192.168.1.34/24}.

Con il prefisso /24 i primi 24 bit identificano la rete e restano 8 bit per gli host. La subnet mask è quindi \textbf{255.255.255.0}.

L'indirizzo di rete si ottiene azzerando tutti i bit host:
\[
192.168.1.34 \rightarrow 192.168.1.0
\]
L'indirizzo di broadcast si ottiene ponendo a 1 tutti i bit host:
\[
192.168.1.255
\]
Il primo host disponibile è \textbf{192.168.1.1}, l'ultimo host disponibile è \textbf{192.168.1.254}.

Gli host utilizzabili sono:
\[
2^8 - 2 = 254
\]

\section*{Esercizio 2}
Rete \textbf{172.16.8.0/23}.

Il prefisso /23 corrisponde alla subnet mask \textbf{255.255.254.0}. Restano 9 bit per la parte host.

Il blocco nel terzo ottetto vale:
\[
256 - 254 = 2
\]
Le reti procedono quindi a passi di 2 nel terzo ottetto. La rete data parte da \textbf{172.16.8.0} e termina prima della successiva, cioè \textbf{172.16.10.0}.

Quindi:
\begin{itemize}
    \item indirizzo di rete: \textbf{172.16.8.0}
    \item indirizzo di broadcast: \textbf{172.16.9.255}
    \item host validi: da \textbf{172.16.8.1} a \textbf{172.16.9.254}
\end{itemize}

Numero massimo di host:
\[
2^9 - 2 = 510
\]

\section*{Esercizio 3}
Dato \textbf{10.0.5.129/26}.

La subnet mask di /26 è \textbf{255.255.255.192}. Restano 6 bit per gli host.

La dimensione del blocco nell'ultimo ottetto è:
\[
256 - 192 = 64
\]
Le subnet dell'ultimo ottetto sono quindi:
\[
0, 64, 128, 192
\]
L'indirizzo 129 appartiene all'intervallo \textbf{128--191}.

Quindi:
\begin{itemize}
    \item rete: \textbf{10.0.5.128}
    \item broadcast: \textbf{10.0.5.191}
    \item primo host: \textbf{10.0.5.129}
    \item ultimo host: \textbf{10.0.5.190}
\end{itemize}

Dimensione del blocco: \textbf{64 indirizzi}. Host utilizzabili:
\[
2^6 - 2 = 62
\]

\section*{Esercizio 4}
Rete iniziale \textbf{192.168.10.0/24}, da suddividere in \textbf{4 subnet uguali}.

Per ottenere 4 subnet servono 2 bit in più per la parte di rete, perché:
\[
2^2 = 4
\]
Il nuovo prefisso è quindi:
\[
/24 + 2 = /26
\]
La nuova subnet mask è \textbf{255.255.255.192}.

Restano 6 bit per gli host, dunque gli host utilizzabili per subnet sono:
\[
2^6 - 2 = 62
\]

La dimensione del blocco è 64, quindi le quattro subnet sono:
\begin{itemize}
    \item \textbf{192.168.10.0/26}
    \item \textbf{192.168.10.64/26}
    \item \textbf{192.168.10.128/26}
    \item \textbf{192.168.10.192/26}
\end{itemize}

\section*{Esercizio 5}
Rete iniziale \textbf{192.168.50.0/24}, da dividere in \textbf{8 subnet}.

Per ottenere 8 subnet servono 3 bit aggiuntivi:
\[
2^3 = 8
\]
Nuovo prefisso:
\[
/24 + 3 = /27
\]
La nuova subnet mask è \textbf{255.255.255.224}.

La dimensione del blocco è:
\[
256 - 224 = 32
\]
Ogni subnet contiene 32 indirizzi totali, di cui 30 utilizzabili.

Le prime tre subnet sono:
\begin{itemize}
    \item \textbf{192.168.50.0/27}
    \item \textbf{192.168.50.32/27}
    \item \textbf{192.168.50.64/27}
\end{itemize}

La terza subnet va da 192.168.50.64 a 192.168.50.95, quindi il broadcast della terza subnet è:
\[
\textbf{192.168.50.95}
\]

\section*{Esercizio 6}
Dato \textbf{172.20.10.75/27}.

La subnet mask di /27 è \textbf{255.255.255.224}. La dimensione del blocco è:
\[
256 - 224 = 32
\]
Le subnet nell'ultimo ottetto sono:
\[
0, 32, 64, 96, 128, 160, 192, 224
\]
L'indirizzo 75 appartiene all'intervallo \textbf{64--95}.

Quindi:
\begin{itemize}
    \item rete: \textbf{172.20.10.64}
    \item broadcast: \textbf{172.20.10.95}
    \item host validi: da \textbf{172.20.10.65} a \textbf{172.20.10.94}
\end{itemize}

\section*{Esercizio 7}
Rete \textbf{10.1.0.0/16} con almeno \textbf{500 host per subnet}.

Occorre trovare il numero minimo di bit host tali che:
\[
2^h - 2 \geq 500
\]
Verifichiamo:
\[
2^8 - 2 = 254 \quad \text{non basta}
\]
\[
2^9 - 2 = 510 \quad \text{basta}
\]
Quindi servono \textbf{9 bit host}.

Il nuovo prefisso è:
\[
32 - 9 = /23
\]
La nuova subnet mask è \textbf{255.255.254.0}.

Poiché la rete iniziale era /16, sono stati presi 7 bit per creare subnet:
\[
23 - 16 = 7
\]
Numero massimo di subnet ottenibili:
\[
2^7 = 128
\]

Host utilizzabili per subnet:
\[
2^9 - 2 = 510
\]

\section*{Esercizio 8}
Dato \textbf{192.168.100.200/28}.

La subnet mask di /28 è \textbf{255.255.255.240}. La dimensione del blocco è:
\[
256 - 240 = 16
\]
Le subnet dell'ultimo ottetto sono:
\[
0, 16, 32, 48, 64, 80, 96, 112, 128, 144, 160, 176, 192, 208, 224, 240
\]
L'indirizzo 200 appartiene all'intervallo \textbf{192--207}.

Quindi:
\begin{itemize}
    \item rete: \textbf{192.168.100.192}
    \item broadcast: \textbf{192.168.100.207}
    \item primo host: \textbf{192.168.100.193}
    \item ultimo host: \textbf{192.168.100.206}
\end{itemize}

\section*{Esercizio 9}
Rete \textbf{192.168.1.0/24}, da suddividere con VLSM per:
\begin{itemize}
    \item reparto A: 100 host
    \item reparto B: 50 host
    \item reparto C: 20 host
\end{itemize}

Nel VLSM si assegna prima la subnet più grande.

\subsection*{Reparto A: 100 host}
Serve il minimo numero di bit host tale che:
\[
2^h - 2 \geq 100
\]
Con 7 bit si ha:
\[
2^7 - 2 = 126
\]
Quindi serve una subnet \textbf{/25} con mask \textbf{255.255.255.128}.

Assegniamo:
\begin{itemize}
    \item rete: \textbf{192.168.1.0/25}
    \item host: da \textbf{192.168.1.1} a \textbf{192.168.1.126}
    \item broadcast: \textbf{192.168.1.127}
\end{itemize}

\subsection*{Reparto B: 50 host}
Serve:
\[
2^6 - 2 = 62
\]
Quindi una subnet \textbf{/26} con mask \textbf{255.255.255.192}.

Lo spazio libero successivo parte da 192.168.1.128:
\begin{itemize}
    \item rete: \textbf{192.168.1.128/26}
    \item host: da \textbf{192.168.1.129} a \textbf{192.168.1.190}
    \item broadcast: \textbf{192.168.1.191}
\end{itemize}

\subsection*{Reparto C: 20 host}
Serve:
\[
2^5 - 2 = 30
\]
Quindi una subnet \textbf{/27} con mask \textbf{255.255.255.224}.

Lo spazio libero successivo parte da 192.168.1.192:
\begin{itemize}
    \item rete: \textbf{192.168.1.192/27}
    \item host: da \textbf{192.168.1.193} a \textbf{192.168.1.222}
    \item broadcast: \textbf{192.168.1.223}
\end{itemize}

\section*{Esercizio 10}
Rete iniziale \textbf{172.16.0.0/22}. Deve essere suddivisa in subnet da almeno \textbf{60 host}.

Cerchiamo i bit host minimi:
\[
2^5 - 2 = 30 \quad \text{non basta}
\]
\[
2^6 - 2 = 62 \quad \text{basta}
\]
Servono quindi 6 bit host, dunque il nuovo prefisso è:
\[
32 - 6 = /26
\]
La nuova subnet mask è \textbf{255.255.255.192}.

La rete iniziale /22 contiene:
\[
2^{32-22} = 2^{10} = 1024 \text{ indirizzi}
\]
Ogni subnet /26 contiene:
\[
2^{32-26} = 2^6 = 64 \text{ indirizzi}
\]
Numero totale di subnet ottenibili:
\[
1024 / 64 = 16
\]

Le prime quattro subnet sono:
\begin{itemize}
    \item \textbf{172.16.0.0/26}
    \item \textbf{172.16.0.64/26}
    \item \textbf{172.16.0.128/26}
    \item \textbf{172.16.0.192/26}
\end{itemize}

La quarta subnet va da 172.16.0.192 a 172.16.0.255, quindi il broadcast della quarta subnet è:
\[
\textbf{172.16.0.255}
\]

\end{document}
